\subsection{Full-map symmetry averaging does not increase the risk}
\label{sec:theory-sym}

The elasticity problem behind the benchmark is invariant under the reflection
$x_1\mapsto 1-x_1$: the domain, the boundary partition, and the constitutive
model are all symmetric, the input distribution has a reflection-invariant
covariance, and reflecting the load reflects the stress field. On the discrete
grid this is expressed by two permutation matrices: $S\in\mathbb R^{p\times p}$
reverses the load samples and $T\in\mathbb R^{q\times q}$ reverses the rows of
the stress field. Both are orthogonal involutions. Supplement~\ref{sec:symcheck}
examines the discrete equivariance relation $G(Su)=T\,G(u)$
empirically. The following proposition assumes this relation exactly.

\begin{proposition}[symmetrization]\label{prop:sym}
Assume $G(Su)=TG(u)$ for $\mu$-a.e.\ $u$, that $\mu$ is invariant under $S$,
and that $T$ is an orthogonal involution ($T^2=I$). For a measurable
$\widehat G:\mathcal U\to\mathcal V$ define its symmetrization
\[
\widehat G_S(u)\;=\;\tfrac12\left(\widehat G(u)+T\,\widehat G(Su)\right).
\]
Then, for every loss $\ell(\cdot,v)$ that is convex in its first argument and
satisfies $\ell(Tw,Tv)=\ell(w,v)$,
\[
\mathbb E_{u\sim\mu}\,\ell\big(\widehat G_S(u),G(u)\big)\;\le\;
\mathbb E_{u\sim\mu}\,\ell\big(\widehat G(u),G(u)\big).
\]
\end{proposition}
\noindent\emph{Proof.} See Supplement~\ref{proof:prop:sym}.


The relative error satisfies both hypotheses when the target norm is
nonzero. The proposition guarantees that test-time reflection averaging
cannot increase population risk under exact equivariance and distributional
invariance. For a finite training sample, Jensen's inequality also shows that
the reflection-augmented empirical loss upper-bounds the empirical loss of
the symmetrized predictor. The objectives are generally different, so the
proposition does not guarantee that training with augmentation improves the
fitted model. The near-reflection comparisons are consistent with symmetry of the
distributed solver output. In the experiment, the averages improve twenty-nine of thirty members across the two MLP variants
and the refiner, and worsen one by a thousandth of a point. For the refiner, equivariance of its conditioning field is a sufficient
additional condition for this guarantee. The convexity
argument is standard; \citet{lyle2020benefits} give a general account.

\paragraph{A supplied conditioning field is an additional input.}
For a fitted refiner $F(u,h)$ and a fixed kernel predictor $h(u)$, define
$f(u)=F(u,h(u))$. The proposition concerns
\[
f_S(u)=\tfrac12\big(F(u,h(u))+T F(Su,h(Su))\big).
\]
The implementation instead uses $T h(u)$ as the second branch's conditioning
field. Thus $h(Su)=T h(u)$ is a sufficient additional condition for identifying
the implemented average with $f_S$. Equivariance of the target and invariance
of the load distribution do not by themselves give this property to a fitted
kernel predictor, and training the refiner with randomly reflected triples
does not force exact equivariance; the non-increase guarantee for its
averaging operation is conditional on that identity.

The effect of this difference can be bounded. Write
$b(u)=\tfrac12(F(u,h(u))+T F(Su,T h(u)))$ for the implemented average and
$R(g)=\mathbb E\|g(u)-G(u)\|_2/\|G(u)\|_2$, assuming nonzero targets
almost surely and finite quantities below. Under the proposition's exact
symmetry assumptions,
\begin{equation}\label{eq:refiner-defect}
\begin{aligned}
R(b)&\le R(f_S)+\frac12\mathbb E
\frac{\|F(Su,T h(u))-F(Su,h(Su))\|_2}{\|G(u)\|_2}
\\
&\le R(f)+\frac12\mathbb E
\frac{\|F(Su,T h(u))-F(Su,h(Su))\|_2}{\|G(u)\|_2}.
\end{aligned}
\end{equation}
Indeed, the triangle inequality bounds the risk difference by
$\mathbb E\|b-f_S\|_2/\|G\|_2$, and orthogonality of $T$ gives the
displayed defect exactly; the second inequality is Proposition~\ref{prop:sym}.
If $F(u,\cdot)$ is uniformly $L_h$-Lipschitz, the additional term is at most
$(L_h/2)\mathbb E[\|T h(u)-h(Su)\|_2/\|G(u)\|_2]$.
This defect and the Lipschitz constant are not measured in the experiments,
so the inequality is not evaluated numerically.

The two averages can have different risks even for an exactly symmetric target. Let the
input take the values $+1$ and $-1$ with equal probability, let $S$ exchange
them, and let $T$ swap two output coordinates. Set $G(u)=(1,1)$,
$h(+1)=(1,1)$ and $h(-1)=(3,3)$. Define $F(+1,z)=z$ and
$F(-1,z)=(4,4)-z$. The four evaluations are
\[
\begin{aligned}
F(+1,h(+1))&=(1,1), & F(-1,T h(+1))&=(3,3),\\
F(-1,h(-1))&=(1,1), & F(+1,T h(-1))&=(3,3).
\end{aligned}
\]
Thus $f(u)=G(u)$ at both inputs, and the full-map average
also has zero error. Averaging with $T h(u)$ instead gives $(2,2)$ at both
inputs, with relative error one. This finite example disproves a general
guarantee for the supplied-field operation and attains equality in
\eqref{eq:refiner-defect} (here $L_h=1$).

\subsection{The kernel-flow regularizer and leave-half-out error}
\label{sec:theory-kf}

The kernel-flow (KF) loss of Owhadi and Yoo \citep{owhadi2019kernel}, in the
$\ell^2$ variant used by \citet{yoo2021deep} to regularize deep networks,
defines the optional regularizer in the ablation study. Let $z_i=\phi_\theta(u_i)$ be the features of a batch
$B=\{1,\dots,b\}$ ($b$ even), let $k$ be a positive definite kernel on feature
space, and for $C\subset B$, $|C|=b/2$, let $I_C$ denote the $k$-interpolant
of $\{(z_j,v_j):j\in C\}$. The batch KF loss is
\begin{equation}\label{eq:kf}
e_2(\theta;C)\;=\;\sum_{i\in B}\big\|v_i-I_C(z_i)\big\|_2^2 .
\end{equation}

\begin{proposition}[KF loss is a leave-half-out estimate]\label{prop:kf}
Let $C$ be drawn uniformly among the subsets of $B$ of size $b/2$ and assume
$k$ restricted to $\{z_i\}_{i\in B}$ is strictly positive definite. Then
$e_2(\theta;C)=\sum_{i\in B\setminus C}\|v_i-I_C(z_i)\|_2^2$, and
\[
\mathbb E_C\big[e_2(\theta;C)\big]\;=\;\frac b2\;
\mathbb E_{C,\,i\,\sim\,\mathrm{Unif}(B\setminus C)}
\big[\big\|v_i-I_C(z_i)\big\|_2^2\big],
\]
i.e.\ $\tfrac2b\,e_2$ is an unbiased estimator, over the split randomness, of
the expected squared error of the kernel predictor on a held-out point
when trained on a random half of the batch. If the parameterized
feature and kernel maps are differentiable and the derivative of
$e_2(\theta;C)$ is bounded in a parameter neighborhood by a random variable
integrable over $(B,C)$, then, when $(B,C)$
are resampled at every step, the stochastic gradient
$\nabla_\theta e_2(\theta;C)$ is unbiased for
$\nabla_\theta\,\mathbb E_{B,C}[e_2(\theta;C)]$.
\end{proposition}
\noindent\emph{Proof.} See Supplement~\ref{proof:prop:kf}.


The implemented ablation uses normalized targets and the mean squared loss
over all $bq$ batch-output entries, with predictor
$K_{BC}(K_{CC}+10^{-3}I)^{-1}Y_C$. Its errors on $C$ need not vanish:
on those rows the residual is
$10^{-3}(K_{CC}+10^{-3}I)^{-1}Y_C$. Consequently the exact leave-half-out
identity above does not apply to the regularized loss. The learned Gaussian
bandwidth multiplier is combined with a batch median distance computed
under \texttt{no\_grad}; the resulting gradient also omits the dependence
of that median on the features. This stop-gradient rule differs from
the full derivative in Proposition~\ref{prop:kf}. Neither the ideal identity nor its regularized
analogue estimates the population risk of the adaptively trained network.


\subsection{Capacity of stacking with fixed members}
\label{sec:theory-stack}

Convexity bounds the loss of a convex combination by the corresponding
weighted average of member losses. A separate uniform-deviation argument
bounds the cost of fitting the weights when the candidate members are fixed
independently of an i.i.d. validation sample and the losses have a population
bound.

\begin{proposition}[validation stacking]\label{prop:stack}
Let $f_1,\dots,f_M:\mathcal U\to\mathcal V$ be fixed maps, let
$\Delta=\{w\in\mathbb R^M:w_m\ge0,\sum_m w_m=1\}$, and for $w\in\Delta$ write
$F_w=\sum_m w_m f_m$. Assume the members' per-sample relative errors are
bounded: $\ell(f_m(u),G(u))\le B$ for $\mu$-a.e.\ $u$ and every $m$.
\begin{enumerate}
\item[(i)] For every $w\in\Delta$,
$\ell(F_w(u),G(u))\le\sum_m w_m\,\ell(f_m(u),G(u))$ pointwise; in particular
$R(F_w)\le\sum_m w_m R(f_m)\le\max_m R(f_m)$, and $\ell(F_w(u),G(u))\le B$.
\item[(ii)] Let $\widehat w$ minimize the empirical risk
$\widehat R(w)=\tfrac1m\sum_{i=1}^m\ell(F_w(\tilde u_i),G(\tilde u_i))$ over
$\Delta$, computed on an i.i.d.\ validation sample
$\tilde u_1,\dots,\tilde u_m\sim\mu$ independent of $f_1,\dots,f_M$. Then for
every $\delta\in(0,1)$, with probability at least $1-\delta$,
\[
R(F_{\widehat w})\;\le\;\min_{w\in\Delta}R(F_w)\;+\;\frac{8B}{m}\;+\;
B\sqrt{\frac{2\big((M-1)\log\!\big(m(M-1)+1\big)+\log(2/\delta)\big)}{m}} .
\]
\end{enumerate}
\end{proposition}
\noindent\emph{Proof.} See Supplement~\ref{proof:prop:stack}.


With $M=5$, $m=1000$ and $\delta=0.05$, the excess term is about
$0.28B$. Substituting $B=0.25$, the largest observed member error, gives
roughly seven percentage points, if that value is a uniform population bound. This excess term is much larger than the measured selection effects. The result supplies the scaling $\sqrt{M\log m/m}$;
the much smaller validation-to-test gaps are empirical measurements.

The stack the pipeline deploys on the structural-mechanics benchmark is more general
than a single convex vector: the weights are affine and vary per grid point,
fitted by ridge on the validation split, and selection between the affine and global convex stacks uses a half-split
comparison: each is fitted on one half, evaluated on the other, and the
one with lower error is refitted on the full split. Proposition~\ref{prop:stack} does not cover that
construction. The following two statements give conditional capacity bounds
for related fixed-candidate selection and affine classes.

\begin{lemma}[validated selection]\label{lem:select}
Let $g_1,\dots,g_K$ be fixed maps from $\mathcal U$ to $\mathcal V$. Assume
$\ell(g_k(u),G(u))\le B$ for $\mu$-a.e.\ $u$ and every $k$, and let
$\widehat k$ minimize the empirical risk over an i.i.d.\ sample of size $m'$
drawn from $\mu$ independently of everything used to build the $g_k$. Then,
with probability at least $1-\delta$,
\[
R(g_{\widehat k})\;\le\;\min_{k\le K}R(g_k)\;+\;
B\sqrt{\frac{2\log(2K/\delta)}{m'}} .
\]
\end{lemma}
\noindent\emph{Proof.} See Supplement~\ref{proof:lem:select}.


The half-split comparison uses $K=2$ and $m'=500$. The lemma applies
only if both candidate stacks were constructed independently of the evaluation
half. Refitting the selected stack changes the predictor. The following
proposition treats affine fitting under independence and boundedness
assumptions.

\begin{proposition}[per-pixel affine stacking]\label{prop:affine}
Fix members $f_1,\dots,f_M$, write $D=\sqrt{M+1}$, and for each output
coordinate $j$ let $\varphi_j(u)=(f_1(u)_j,\dots,f_M(u)_j,b)\in\mathbb R^{M+1}$,
where $b$ bounds all member and target values:
$|G(u)_j|\le b$, $|f_m(u)_j|\le b$ for $\mu$-a.e.\ $u$ and all $j,m$. Consider
the per-coordinate affine class
$\{u\mapsto w^\top\varphi_j(u):\|w\|_2\le W\}$ and the coordinate risks
$R_j(w)=\mathbb E\,(w^\top\varphi_j(u)-G(u)_j)^2$, estimated on an i.i.d.\
validation sample of size $m$ independent of the members.
\begin{enumerate}
\item[(i)] If $\widehat w_j$ minimizes the empirical risk of coordinate $j$
over the class, then with probability at least $1-\delta$,
\[
\begin{aligned}
&\frac1q\sum_{j=1}^q\Big[R_j(\widehat w_j)-\min_{\|w\|_2\le W}R_j(w)\Big]\\
&\qquad\le\frac{8\,W b^2 D\,(1+WD)+2\,b^2(1+WD)^2\sqrt{\tfrac12\log(4q/\delta)}}{\sqrt m}.
\end{aligned}
\]
\item[(ii)] If instead $\widehat w_j$ minimizes the ridge objective
$\widehat R_j(w)+\lambda\|w\|_2^2$ with $\lambda\ge0$ and satisfies $\|\widehat w_j\|_2\le W$,
the same bound holds with an additional $\lambda W^2$ on the right.
\end{enumerate}
\end{proposition}
\noindent\emph{Proof.} See Supplement~\ref{proof:prop:affine}.


For fixed $M,b,W$, this is an $m^{-1/2}$ bound on coordinatewise
mean-squared risk, with logarithmic dependence on the number of coordinates.
It is not a relative-error bound. Its constants depend polynomially on $b$
and $W$; substituting the measured values makes the bound exceed the realized
excess by many orders of magnitude. Those substitutions are illustrative,
since a finite observed maximum is not a population bound.

A fixed grid of residual corrections gives a further composite class. For
each configuration, correction of training residuals is affine in the stack
weights, so the same uniform-deviation argument applies with a transformed
feature map. The following result requires the training arrays, members and
correction configurations to be fixed independently of validation.

\begin{corollary}[capacity bound for a fixed composite correction class]
\label{cor:certificate}
Assume the setting of Proposition~\ref{prop:affine}, and let $\Gamma$ be a
finite set of correction configurations, each defining weights
$c_\gamma(u)\in\mathbb R^{n_{\mathrm{tr}}}$ with the corrected predictor
\[
h_{w,\gamma}(u)_j\;=\;w_j^\top\varphi_j(u)
+c_\gamma(u)^\top\big(Y_j-\Phi_j w_j\big),
\]
where $Y_j\in\mathbb R^{n_{\mathrm{tr}}}$ collects the training targets of
coordinate $j$, and $\Phi_j\in\mathbb R^{n_{\mathrm{tr}}\times(M+1)}$
has rows $\varphi_j(u_i)^\top$, including the final column of intercept
features equal to $b$; $\gamma=0$
(no correction, $c_0\equiv0$) is included in $\Gamma$. Suppose the correction
weights are uniformly bounded, $\sup_{u,\gamma}\|c_\gamma(u)\|_1\le\kappa$.
Assume that $Y_j$, $\Phi_j$, the members and all $c_\gamma$ are fixed
independently of the validation sample, and that their training target and
feature entries satisfy the same absolute bound $b$. For any
validation-measurable choice
$(\widehat w,\widehat\gamma)$ with $\|\widehat w_j\|_2\le W$ and
$\widehat\gamma$ minimizing the empirical validation risk, aggregated over
coordinates, among $\gamma\in\Gamma$ at the fitted $\widehat w$ (the grid
choice is shared across coordinates), with
probability at least $1-\delta$,
\[
\begin{aligned}
\frac1q\sum_j R_j\big(h_{\widehat w,\widehat\gamma}\big)
&\le \min_{\gamma\in\Gamma}\frac1q\sum_j R_j\big(h_{\widehat w,\gamma}\big)\\
&\quad+\frac{2(1+\kappa)^2 b^2(1+WD)}{\sqrt m}
\Bigg(4WD+(1+WD)\sqrt{\frac{\log(4q|\Gamma|/\delta)}{2}}\Bigg).
\end{aligned}
\]
Since $\gamma=0$ is among the candidates, the oracle on the right is at most
the fitted stack's risk, and combining with
Proposition~\ref{prop:affine}(ii) bounds the composite risk against the best
affine stack, provided the hypotheses of both results hold (with failure
probabilities allocated between them). Here $b$ and $W$ must be population
and class bounds; the observed values are
diagnostics of them. For
kernel ridge weights, $\kappa$ is supplied by the next lemma.
\end{corollary}
\noindent\emph{Proof.} See Supplement~\ref{proof:cor:certificate}.


For kernel ridge regression, the correction-weight bound follows analytically. In
the pipeline each $\gamma=(s,\lambda)$ is a length scale and a nugget, and
the weights are those of kernel ridge regression,
$c_\gamma(u)=(K_\gamma+n\lambda I)^{-1}k_\gamma(X,u)$, so that
$\widehat r(u)=c_\gamma(u)^\top R$ in the notation of
Section~\ref{sec:method-kernel}. Their $\ell^1$ norm is bounded uniformly in
$u$ by the nugget alone.

\begin{lemma}[uniform bound on kernel ridge weights]\label{lem:kappa}
Let $k$ be a positive definite kernel with $k(u,u)\le k_{\max}$ for all $u$,
let $X=(u_1,\dots,u_n)$ be any design with Gram matrix $K$, and let
$\lambda>0$. The kernel ridge weights
$c_\lambda(u)=(K+n\lambda I)^{-1}k(X,u)$ satisfy, for every $u$,
\[
\|c_\lambda(u)\|_1\;\le\;\sqrt n\,\|c_\lambda(u)\|_2\;\le\;
\frac12\sqrt{\frac{k(u,u)}{\lambda}}\;\le\;\frac12\sqrt{\frac{k_{\max}}{\lambda}} ,
\]
and the constant $\tfrac12$ cannot be improved: for every $\lambda\in(0,1)$
there are a kernel, a two-point design and a point $u$ at which the first
and last members are equal. In particular, for the residual correction of
Section~\ref{sec:method-kernel} --- a Mat\'ern kernel with $k(u,u)=1$ and
the nugget grid $\lambda\in\{10^{-6},10^{-5},10^{-3}\}$ --- the hypothesis
of Corollary~\ref{cor:certificate} holds, whatever the length scale and the
design, with
\[
\kappa_{\mathrm{an}}\;:=\;\tfrac12\max_{\gamma\in\Gamma}\lambda_\gamma^{-1/2}\;=\;500.
\]
\end{lemma}
\noindent\emph{Proof.} See Supplement~\ref{proof:lem:kappa}.


Positive semidefiniteness of the Gram matrix of the design together with
$u$ confines
$k(X,u)$ to the range of $K$ with $k(X,u)^\top K^{+}k(X,u)\le k(u,u)$, and
$\mu/(\mu+n\lambda)^2\le1/(4n\lambda)$ for every eigenvalue $\mu$.
Thus $\kappa_{\mathrm{an}}$ is a uniform bound determined by the kernel
diagonal and nugget grid. The excess term contains $(1+\kappa)^2$;
with $\kappa_{\mathrm{an}}=500$, it exceeds the measured errors by many
orders of magnitude. The realized weights have a smaller scale:
on a test probe of five hundred points the largest $\ell^1$ norm over the
grid averages $\kappaTen$ over the ten seeds
(Section~\ref{app:verification}), about a fifteenth of
$\kappa_{\mathrm{an}}$. Substituting that value would reduce the excess term
by about a factor of 240, but a finite-probe maximum does not supply the
required uniform bound.

The corollary bounds coordinatewise mean-squared risk for candidates fixed
independently of validation. The deployed estimator instead selects by
mean relative $L^2$ error, tunes its kernel stages on an $8000$-row
subsample, refits the selected kernel on all $19000$ rows, and reuses the
validation split for checkpoint and member selection. These operations are
outside the corollary's assumptions. Likewise, observed maxima do not
establish the population bounds $B,b,L,W$ required by the preceding
capacity results.
